% from ega2-1-fr.tex
\oldpage[II]{83}

Ce lemme étant admis, il existe une immersion fermée $j_1$ de $P$ dans
\[
  P_1=\mathbf{P}(g^*(\mathcal{F})\otimes\mathcal{K}^{\otimes(-m)})
\]
telle que $\mathcal{O}_P(1)$ soit isomorphe à
$j_1^*(\mathcal{O}_{P_1}(1))$ (\hyperref[II.4.1.2-fr]{4.1.2}). D'autre part,
il existe un isomorphisme de $P_1$ sur
$P_2=\mathbf{P}(g^*(\mathcal{F}))$, identifiant
$\mathcal{O}_{P_1}(1)$ à
$\mathcal{O}_{P_2}(1)\otimes_Y\mathcal{K}^{\otimes(-m)}$
(\hyperref[II.4.1.4-fr]{4.1.4})~; on a donc une immersion fermée
$j_2:P\to P_2$ telle que $\mathcal{O}_P(1)$ soit isomorphe à
$j_2^*(\mathcal{O}_{P_2}(1))\otimes_Y\mathcal{K}^{\otimes(-m)}$. Enfin,
$P_2$ s'identifie à $P_3\times_ZY$, où $P_3=\mathbf{P}(\mathcal{F})$, et
$\mathcal{O}_{P_2}(1)$ à
$\mathcal{O}_{P_3}(1)\otimes_Z\mathcal{O}_Y$
(\hyperref[II.4.1.3-fr]{4.1.3}). Par définition,
$\mathcal{O}_{P_3}(1)$ est très ample pour $Z$~; comme il en est de même de
$\mathcal{K}$, on conclut de (iv) que
$\mathcal{O}_{P_2}(1)\otimes_Y\mathcal{K}$ est très ample pour $Z$~; il en
est de même de
$\mathcal{M}=j_2^*(\mathcal{O}_{P_2}(1)\otimes_Y\mathcal{K})$ en vertu de
(i bis), et $\mathcal{O}_P(1)$ est isomorphe à
$\mathcal{M}\otimes_Y\mathcal{K}^{\otimes(-m-1)}$, ce qui achève la
démonstration.
\end{proof}

\begin{proposition}[4.4.11]
\label{II.4.4.11-fr}
Soient $f:X\to Y$, $f':X'\to Y$ deux morphismes, $X''$ le préschéma somme
$X\sqcup X'$, $f''$ le morphisme $X''\to Y$ qui coïncide avec $f$ (resp.
$f'$) dans $X$ (resp. $X'$). Soit $\mathcal{L}$ (resp. $\mathcal{L}'$) un
$\mathcal{O}_X$-Module (resp. un $\mathcal{O}_{X'}$-Module) inversible, et
soit $\mathcal{L}''$ le $\mathcal{O}_{X''}$-Module inversible qui coïncide
avec $\mathcal{L}$ dans $X$ et avec $\mathcal{L}'$ dans $X'$. Pour que
$\mathcal{L}''$ soit très ample relativement à $f''$, il faut et il suffit
que $\mathcal{L}$ soit très ample relativement à $f$ et $\mathcal{L}'$ très
ample relativement à $f'$.
\end{proposition}

\begin{proof}
On est aussitôt ramené au cas où $Y$ est affine. Si $\mathcal{L}''$ est très
ample il en est de même de $\mathcal{L}$ et $\mathcal{L}'$ en vertu de
(\hyperref[II.4.4.10-fr]{4.4.10, (i bis)}). Inversement, si $\mathcal{L}$ et
$\mathcal{L}'$ sont très amples, il résulte aussitôt de la définition
(\hyperref[II.4.4.2-fr]{4.4.2}) et de
(\hyperref[II.4.3.6-fr]{4.3.6}) que $\mathcal{L}''$ est très ample.
\end{proof}

\subsection{Faisceaux amples}
\label{subsection:II.4.5-fr}

\begin{env}[4.5.1]
\label{II.4.5.1-fr}
Étant donnés un préschéma $X$ et un $\mathcal{O}_X$-Module inversible
$\mathcal{L}$, nous poserons, pour tout $\mathcal{O}_X$-Module $\mathcal{F}$
(lorsque aucune confusion ne sera possible sur $\mathcal{L}$),
$\mathcal{F}(n)=\mathcal{F}\otimes_{\mathcal{O}_X}\mathcal{L}^{\otimes n}$
($n\in\mathbf{Z}$)~; nous poserons d'autre part
\[
  S=\bigoplus_{n\geq0}\Gamma(X,\mathcal{L}^{\otimes n})
\]
(sous-anneau gradué de l'anneau $\Gamma_\bullet(\mathcal{L})$ défini dans
(\hyperref[0.5.4.6-fr]{0, 5.4.6})). Si on considère $X$ comme un
$\mathbf{Z}$-préschéma et qu'on désigne par $p$ le morphisme structural
$X\to\operatorname{Spec}(\mathbf{Z})$, il y a correspondance biunivoque
entre les homomorphismes
\[
  p^*(\widetilde{S})\longrightarrow
  \bigoplus_{n\geq0}\mathcal{L}^{\otimes n}
\]
de $\mathcal{O}_X$-Algèbres graduées et les endomorphismes de l'anneau gradué
$S$ (\hyperref[I.2.2.5-fr]{I, 2.2.5})~; l'homomorphisme
\[
  \varepsilon:p^*(\widetilde{S})\longrightarrow
  \bigoplus_{n\geq0}\mathcal{L}^{\otimes n}
\]
qui correspond ainsi à l'automorphisme identique de $S$ sera dit canonique.
Il lui correspond (\hyperref[II.3.7.1-fr]{3.7.1}) un morphisme
$G(\varepsilon)\to\operatorname{Proj}(S)$ qui sera aussi dit canonique.
\end{env}

\begin{theorem}[4.5.2]
\label{II.4.5.2-fr}
Soient $X$ un schéma quasi-compact ou un préschéma dont l'espace sous-jacent
est noethérien, $\mathcal{L}$ un $\mathcal{O}_X$-Module inversible, $S$
l'anneau gradué
$\bigoplus_{n\geq0}\Gamma(X,\mathcal{L}^{\otimes n})$. Les conditions
suivantes sont équivalentes~:
\begin{enumerate}
  \item[\rm{a)}] Lorsque $f$ parcourt l'ensemble des éléments homogènes de
  $S_+$, les $X_f$ forment une base de la topologie de $X$.
  \item[\rm{a')}] Lorsque $f$ parcourt l'ensemble des éléments homogènes de
  $S_+$, ceux des $X_f$ qui sont affines forment un recouvrement de $X$.
  \item[\rm{b)}] Le morphisme canonique
  $G(\varepsilon)\to\operatorname{Proj}(S)$
  (\hyperref[II.4.5.1-fr]{4.5.1}) est partout défini et est une immersion
  ouverte dominante.

